% Generated draft: requires mathematical and source review.
\documentclass[11pt]{article}
\usepackage[margin=25mm]{geometry}
\usepackage{fontspec}
\setmainfont{Latin Modern Roman}
\usepackage{amsmath,amssymb,mathtools,unicode-math}
\setmathfont{Latin Modern Math}
\usepackage{enumitem,needspace,xurl,hyperref}
\setlength{\parindent}{0pt}
\setlength{\parskip}{5pt}
\setlength{\emergencystretch}{3em}
\newcommand{\N}{\mathbb{N}}
\newcommand{\Z}{\mathbb{Z}}
\newcommand{\Q}{\mathbb{Q}}
\newcommand{\R}{\mathbb{R}}
\newcommand{\C}{\mathbb{C}}
\newcommand{\F}{\mathbb{F}}
\newcommand{\A}{\mathbb{A}}
\newcommand{\PP}{\mathbb{P}}
\newcommand{\E}{\mathbb{E}}
\newcommand{\eps}{\varepsilon}
\newcommand{\dd}{\,\mathrm{d}}
\begin{document}
\section*{UnsolvedMath ID: 20002363}
\textbf{Source problem code:} catalogue entry
\subsection{Definitions and mathematical statement}
The recovered field is
\[
 F=\operatorname{Frac}\bigl(\R[x,y]/(x^2+y^2+1)\bigr),
\]
not a fraction field over the sought valuation ring. This distinction between an ordinary italic ring symbol and blackboard-bold \(\R\) was lost in the registry extraction. A valuation ring \(\mathcal O\subseteq F\) is a subring such that for each \(z\in F^\times\), either \(z\in\mathcal O\) or \(z^{-1}\in\mathcal O\). Nontrivial means \(\mathcal O\ne F\). Moret-Bailly's Question 26 asks whether a nontrivial valuation ring of \(F\) is definable, and asks the analogous question for a proper nonempty finite intersection of valuation rings, asserting that the two existence questions are equivalent. Definability is read in the field ring language; whether parameters from all of \(F\), only \(\R\), or none are allowed is a material convention. The aid treats arbitrary \(F\)-parameters as the unresolved version.

The supplied reduction concerns \emph{geometric} valuations, namely nontrivial valuations trivial on \(\R^\times\). Let \(C/\R\) be the smooth projective conic \(X^2+Y^2+Z^2=0\), with function field \(F\). Its geometric valuation rings are the local discrete valuation rings \(\mathcal O_P\) at closed points \(P\) of \(C\). Choose distinct closed points \(P_1,\ldots,P_r\), \(r\ge1\), and put
\[
 A=\bigcap_{j=1}^r\mathcal O_{P_j}.
\]
For a chosen \(i\), weak approximation supplies \(\pi_i\in F\) such that
\[
 v_{P_i}(\pi_i)>0,\qquad v_{P_j}(\pi_i)=0\quad(j\ne i),
\]
so \(\pi_i\in A\). The aid's positive-existential extraction formula is
\[
 x\in\mathcal O_{P_i}\quad\Longleftrightarrow\quad
 \exists s,u,w\in A\quad(sx\in A\ \wedge\ us+w\pi_i=1).
\]
The equation makes \(s\) a unit at \(P_i\); at the other chosen points the parameter \(\pi_i\) is a unit. This formula is over the predicate for \(A\). If \(A\) is definable in \(F\), substitution gives a definition of \(\mathcal O_{P_i}\), with the extra field parameter \(\pi_i\). If \(A\) is existentially definable, the resulting definition is also existential. Conversely a geometric valuation ring is a one-member intersection. Thus the two existence questions are equivalent for geometric valuations when these additional field parameters are permitted, even at the existential level. This argument is not silently extended to arbitrary nongeometric valuation rings.

The aid also reports a parameter obstruction. The action of \(\operatorname{SO}_3(\R)\) on \(C\) induces \(\R\)-automorphisms of \(F\) and is transitive on its infinite set of closed points. A real-parameter definition must be invariant under every such automorphism. Consequently no individual geometric valuation ring \(\mathcal O_P\), and no intersection of a nonempty finite set of these rings, is definable using only real parameters. The associated finite set of geometric valuation overrings would have to be invariant under this transitive action. This obstruction addresses only the stated geometric class and only real-parameter definitions. It does not rule out definitions using nonconstant elements of \(F\), nor settle the source's unrestricted definability question.
\subsection{Short English statement}
The field is the function field of the anisotropic real conic. For geometric valuations, a definable finite intersection yields a chosen valuation ring by an existential Bezout formula with one extra parameter. The aid rules out real-parameter definitions in this geometric class, while leaving arbitrary field-parameter definability unresolved.
\subsection{Sources}
\begin{itemize}
\item Local registry, permanent UnsolvedMath ID \(20002363\), exact public code \texttt{catalogue entry}; canonical location \texttt{aim-logic-notes.json}, \texttt{notes[138]}.
\item AIM, \emph{Extensions of Hilbert's tenth problem}, Question 26 attributed to Moret-Bailly; supplied reference \url{https://aimath.org/WWN/hilberts10th/hilberts10th.pdf}. No proposer year, original remarks, or literature field is supplied.
\item Supplied verification history reports a prior download and inspection of the page-6 content stream: the sought ring symbol used font \texttt{CMMI12}, whereas the real constant field used \texttt{MSBM10}. This is the record's basis for recovering \(\R\) in the field definition; no PDF was downloaded or inspected in preparing this draft.
\item Supplied research attempt 1: \texttt{reduction}, run status \texttt{open}, rigor label \texttt{proved}, low novelty confidence; registry status \texttt{partially\_solved}, classification \texttt{PARTIAL-PROGRESS}, created and updated 14 August 2026. No primary citation for the geometric extraction or symmetry obstruction is supplied. The aid is machine generated, not peer reviewed, and requires independent verification before citation.
\end{itemize}
\end{document}
